\chapter{Introduction}
\label{chap:introduction}

Following the integration of the linear motor, the unicycle robot contains the main actuation components required for both longitudinal motion and lateral balancing. The wheel motor controls the longitudinal motion of the robot, while the FAULHABER LM2070 linear motor moves a mass along the lateral rod to generate the moment required for roll stabilization. In principle, this configuration provides a complete platform capable of maintaining balance and moving forward at a desired speed. In practice, however, achieving reliable lateral balance remains considerably more difficult than suggested by the idealized model.

Unlike a reaction wheel, the lateral moving mass does not directly provide an arbitrary and continuously available body torque. Its balancing effect depends on the rod position, velocity, and acceleration, while its motion is restricted by the available motor force and finite rod travel. The controller must therefore move the mass quickly enough to prevent the lean angle from growing, but must also stop or reverse the mass before it reaches the mechanical end stops. This creates an important conflict between rapid disturbance rejection and limited actuator travel. End-stop collisions, unmodelled rolling dynamics, sensor error, measurement latency, and discrete-time implementation further increase the difference between theoretical simulations and the physical system.

The objective of this report is not to present a final balancing solution, but to document the modelling, controller-development, simulation, and experimental work performed after the installation of the linear motor. Several control approaches are investigated, including LQR, pole placement, hierarchical control, one-step force optimization, and a preliminary model-predictive-control formulation. Although the reduced model is mathematically controllable, the simulations show that the feasible controller design is strongly constrained by the LM2070 force and travel limits, as well as by the accuracy and latency of the available state measurements.

The report also examines the underlying system dynamics to explain why the problem is difficult. The equilibrium curve, open- and closed-loop eigenvectors, pole clustering, and the influence of physical parameters are analyzed to identify the dominant limitations of the present mechanism. These observations motivate two directions for future work: increasing the lateral moving mass and investigating a hybrid actuator that combines the sustained gravitational moment of the moving mass with the fast transient torque of a reaction wheel.

Chapter~\ref{chap:modelling} introduces the reduced lateral model and discusses the rolling and collision effects that are not fully represented by it. Chapter~\ref{chap:hardware} describes the physical hardware, sensing system, embedded implementation, and experimentally identified parameters. Chapter~\ref{chap:controller-simulation} presents the controller designs and numerical simulations. The subsequent chapter examines the equilibrium geometry, modal structure, and effects of the mechanical parameters. Finally, Chapter~\ref{chap:further-works} summarizes the proposed hardware modifications and future control directions.